\documentclass[11pt]{article}
\usepackage[T1]{fontenc}
\usepackage{lmodern,amsmath,amssymb,booktabs}
\usepackage[margin=1in]{geometry}
\usepackage[hidelinks]{hyperref}
\title{How many forecasts preserve the benefit of routing financial news?}
\author{Lluis Eriksson}
\date{20 November 2026}
\begin{document}
\maketitle
\begin{abstract}
Six forecasts produced from different pairs of news articles can lose their accuracy advantage when only two outputs are retained. We reanalyse published results for 609 resolved prediction-market contracts. We derive the expected Brier loss for every way of distributing a retained panel across three evidence groups, with uniform sampling within each configuration. Against the complete six-output shared-news baseline, retaining at least three routed forecasts suffices for all three studied model endpoints. The retained number and group counts can even vary across contracts. Three is the smallest such guarantee for GPT and OLMo. For Qwen, one uniformly selected forecast suffices. With two outputs, choosing distinct groups favors routing on all endpoints, whereas choosing both from the same group reverses the comparison for GPT and OLMo. These are finite-benchmark expectations over output selection. Replication on the original forecasts confirms the sample thresholds. Paired event-bootstrap intervals across contracts include zero for all three endpoints.
\end{abstract}

\section{A decision after the forecasts are available}
Aluguvelly [1] compared six-member forecasting panels for binary outcomes. In the shared condition (C1) every member received the same two news articles. In the routed condition (C3) six members were divided evenly among the three pairs of three articles. The number of article exposures was twelve in each condition. On the 609 contracts observed before market end, the mean six-member Brier loss favored C3 descriptively at each of three endpoints. Only the Qwen net improvement survived the original within-endpoint multiplicity correction. A contract's Brier loss for forecast $p$ and outcome $y$ is $(p-y)^2$.

Suppose an aggregator has these six outputs and can retain only some of them. We seek a minimum number that preserves the direction of the empirical mean comparison, allowing the aggregator to choose the number taken from each evidence group. Selecting outputs can also stand in for a particular model of availability. The underlying study used all six calls; our calculations do not demonstrate reduced inference cost for a sequential deployment.
\newpage
\section{A finite-panel identity}
Fix one six-output panel and its outcome. Index its three pairs by $g$. Let $\mu_g$ be a pair mean, $u$ the six-member mean, $B$ the average of $(\mu_g-u)^2$ over the three groups and $W$ the average of $(p_i-\mu_{g(i)})^2$ over the six members. The full-panel loss is $E=(u-y)^2$. Total dispersion $D=B+W$ equals the average member loss minus $E$, the standard squared-error ensemble decomposition [2].

A configuration $\ell=(\ell_1,\ell_2,\ell_3)$ has entries in $\{0,1,2\}$ and sum $k>0$. Assign its entries to the named groups by a uniform random permutation, then draw each group's requested number of distinct members uniformly, independently between groups. For C1, average over the uniform partition into three artificial pairs, as in [1]. The actual article identities are held fixed in C3.

\paragraph{Proposition.} For the arithmetic mean $\hat p$ of the retained outputs,
\[
 \mathbb E[(\hat p-y)^2\mid p,y,\ell]=E+\alpha_\ell B+\beta_\ell W,\qquad
 \alpha_\ell=\frac{3\sum_j\ell_j^2-k^2}{2k^2},\quad
 \beta_\ell=\frac{\#\{j:\ell_j=1\}}{k^2}.
\]
To see this, write the selected mean as a weighted average of group means plus centered within-group errors. A uniformly permuted centered triple has variance $B$ per entry and covariance $-B/2$ for distinct entries since the triple sums to zero. Expanding the variance with weights $\ell_j/k$ gives $\alpha_\ell B$. Conditional on the group assignment, only singleton draws introduce within-group variance. For one such draw its contribution is the group's mean squared deviation divided by $k^2$. Averaging the assignment gives $\beta_\ell W$. The conditional error has zero mean, so the cross term vanishes. The retained mean has expectation $u$.

\begin{center}
\begin{tabular}{crr@{\qquad}crr}
\toprule
Counts&$\alpha$&$\beta$& Counts&$\alpha$&$\beta$\\\midrule
100&1&1&211&1/16&1/8\\
110&1/4&1/2&220&1/4&0\\
200&1&0&221&1/25&1/25\\
111&0&1/3&222&0&0\\
210&1/3&1/9&&&\\\bottomrule
\end{tabular}
\end{center}
The digits are counts in three unnamed groups. For example, 210 retains all outputs from one group, one from a second, and none from the third.
\newpage
\section{A guarantee against the complete shared panel}
Allow the choice of configuration and panel size to depend arbitrarily on the contract, even on its six routed forecasts. After that choice the groups and members must still be selected by the uniform rule above. For $k\ge3$, both coefficients are at most $1/3$. Consequently each contract's expected routed loss is at most $E_3+D_3/3$. Subtracting the complete shared-panel loss $E_1$ and averaging gives
\[
  \overline{\mathbb E[L_{\rm routed}]-E_1}
  \ \le\ \overline{E_3-E_1}+\overline{D_3}/3 .
\]
This pointwise argument permits contract-dependent choices that would invalidate a direct substitution of a common configuration into mean component changes.

\begin{center}
\begin{tabular}{lrrr}
\toprule
Endpoint& $\overline{E_3-E_1}$ & $\overline{D_3}$ & Upper bound, $k\ge3$\\
\midrule
GPT-4.1-mini & $-0.00486$&0.0119&$-0.000871$\\
Qwen3-30B & $-0.01101$&0.0091&$-0.007955$\\
OLMo-3.1-32B & $-0.00418$&0.0113&$-0.000391$\\
\bottomrule
\end{tabular}
\end{center}
The inputs come from Tables 2 and 3 of [1]. For an upper bound we first add half a unit in the last displayed place to each input, then round the resulting upper bound upward. The guarantee compares a thinned C3 panel to \emph{all six} C1 forecasts. It concerns the mean over the 609 contracts. Individual losses may increase.

For every nonempty configuration the two coefficients are at most one. For Qwen the upper endpoint of $\overline{E_3-E_1}+\overline{D_3}$ is $-0.001855$. Even a single uniformly chosen routed output therefore gives a lower mean loss than all six shared outputs. This conclusion is specific to this benchmark and its Qwen endpoint.

\section{Why two do not suffice for GPT and OLMo}
Now thin both conditions to two outputs by the same rule. Write $\Delta$ for the C3 minus C1 mean component difference. Conditional on two from one group the contrast is $s=\Delta E+\Delta B$. For one from each of two groups it is $t=\Delta E+\Delta B/4+\Delta W/2$. Choose the same group with a fixed probability $q$; the contrast becomes $qs+(1-q)t$. The C1 pseudo-partition makes its marginal pair uniform for every $q$. Its expected loss is at least its full-panel loss. Thus a positive contrast against thinned C1 also supplies a counterexample against complete C1.
\newpage
\begin{center}
\begin{tabular}{lrrr}
\toprule
Endpoint& Same group $s$ & Distinct groups $t$ & Zero crossing $q$\\
\midrule
GPT&$[.004185,.004295]$&$[-.002528,-.002442]$&$[.3625,.3766]$\\
Qwen&$[-.004665,-.004555]$&$[-.009453,-.009367]$&none in $[0,1]$\\
OLMo&$[.003065,.003175]$&$[-.002348,-.002262]$&$[.4160,.4338]$\\
\bottomrule
\end{tabular}
\end{center}
The inputs for $\Delta B,\Delta W$ are the changes in Table 10 of [1] after averaging the C1 pseudo-partitions. The enclosures propagate rounding errors only. The zero crossing is $-t/(s-t)$ and the denominator is positive over each endpoint's rounding box. On a coordinate segment the ratio is monotone or constant, so evaluating the eight vertices bounds its range. Uniformly retaining any two different calls gives $q=1/5$. Setting $q=1$ instead reverses the sign for GPT and OLMo. Combined with the preceding guarantee, three is their smallest count that works uniformly for all allowed policies retaining at least that count.

\paragraph{Interpretation and reproducibility.}
The result compares retrospective policies applied to fixed forecasts. A failure mechanism that preferentially drops a named article group, or selects a member using its probability, can violate conditional symmetry. If forecasts are rerun, decoding variation and changes in evidence may also change the answer. No market prices or trading returns are used. Statistical tests reported for six-output ensembles do not transfer automatically to these policies. The individual-contract replication on the next page quantifies uncertainty for the retention analysis.

The public package contains the transcribed table entries, exact rational rounding bounds, both enumeration scripts and an independent check of the coefficient table using subset sums. The exhaustive code checks all occupied shapes for ternary panels and both outcomes. The preceding variance proof applies to arbitrary real forecasts. No extra model calls are required.

\newpage
\section{Replication on individual contracts}
After the aggregate calculation, we retrieved Aluguvelly's original forecast archive through the annotation in the ACL PDF [1]. Its readme and the author's code include a uniform-random-subset analysis by ensemble size. Our comparison conditions on evidence-group occupancy and permits a different worst-case occupancy in every contract. These changes determine which policies the bounds protect against. We matched the pre-end manifest to the six shared and six routed calls per endpoint, verified the observed article-pair group sizes, and recovered the full-panel contrasts in Table 3 of [1] to its displayed precision.

For each market and each allowed configuration we averaged the squared errors over \emph{all} eligible subsets and subtracted the loss of the complete C1 panel. We took the largest of these differences for each market before computing their arithmetic mean. This maximization uses outcome information and provides an attainable mathematical upper envelope for the symmetric policies; it is not a proposed forecasting strategy. The means below use the minimum counts derived on page 3.

\begin{center}
\begin{tabular}{lrrrr}
\toprule
Endpoint & Minimum & Mean worst contrast & Worse contracts & Paired 95\% interval\\\midrule
GPT &3&$-.001233$&324/609&$[-.008742,.006295]$\\
Qwen &1&$-.001927$&281/609&$[-.009009,.004910]$\\
OLMo &3&$-.000836$&318/609&$[-.007873,.006112]$\\
\bottomrule
\end{tabular}
\end{center}
All three benchmark means favor routing, yet the worst permitted configuration raises loss on more than half of contracts for GPT and OLMo. With two same-group outputs the mean loss difference against full C1 is positive for GPT ($+.005275$) and OLMo ($+.004661$). For Qwen it remains negative ($-.003511$).

We drew 10,000 paired bootstrap samples of parent events, retaining all conditions and endpoints within each event. The 609 strict contracts have 609 distinct parent events. The displayed two-sided percentile intervals include zero. One-sided percentile upper limits with a Bonferroni adjustment across these three contrasts are $+.006883$, $+.005552$ and $+.006762$, respectively (rounded upward). These exploratory intervals depend on exchangeability of events and do not correct for all choices involved in designing this secondary analysis. They do not establish a population-wide improvement. Their sampling interpretation is separate from rounding enclosures for the fixed benchmark.

\paragraph{Model assessments and review context.}
An OpenAI model participated in the analysis and writing. The package associates the first assessment of the exact PDF with its hash and declares this involvement. No independent human review was obtained.

\paragraph{References.}
[1] A. G. Aluguvelly (2026). What Ensemble Brier Scores Hide: Evidence Routing at a Fixed Article-Exposure Budget in LLM Prediction-Market Forecasting. \emph{Proceedings of FinNLP}, pp.\ 373--382. \href{https://aclanthology.org/2026.finnlp-1.30/}{ACL Anthology 2026.finnlp-1.30}.

\smallskip\noindent
[2] A. Krogh and J. Vedelsby (1994). Neural Network Ensembles, Cross Validation, and Active Learning. \emph{Advances in Neural Information Processing Systems}, 7. \href{https://papers.nips.cc/paper_files/paper/1994/hash/b8c37e33defde51cf91e1e03e51657da-Abstract.html}{Proceedings}.
\end{document}